%========================================================================================
%  CAPITOLO 1 — ONDE GRAVITAZIONALI
%  Obiettivo: arrivare, partendo dalla relatività generale, al fondo stocastico di onde
%  gravitazionali (GWB) e al suo spettro a legge di potenza con gamma = 13/3, che è
%  il segnale che useremo nelle iniezioni dei capitoli 5-6.
%  Lunghezza indicativa: 5-6 pagine.
%  Fonti principali: Maggiore vol. 1 e 2, Taylor (2021) cap. 2 e 4, Shaifullah (2025) sez. 4.
%========================================================================================

\chapter{Onde gravitazionali} \label{chpt:GW}
\sloppy

Le onde gravitazionali sono perturbazioni dello spaziotempo che si propagano alla velocità della luce, generate da masse in moto accelerato. Previste da Einstein già nel 1916 come soluzioni ondulatorie delle equazioni di campo linearizzate~\cite{Einstein1916, Einstein1918}, sono state osservate per la prima volta in modo indiretto attraverso il decadimento orbitale della pulsar binaria PSR B1913+16~\cite{HulseTaylor1975, TaylorWeisberg1982} e, un secolo dopo la loro previsione, in modo diretto dagli interferometri LIGO e Virgo, con la rivelazione del segnale GW150914 emesso dalla coalescenza di due buchi neri di massa stellare~\cite{LIGO2016}.

In questo capitolo richiamo gli elementi di teoria necessari per il resto della tesi: la descrizione delle onde gravitazionali in relatività generale linearizzata (Sezione~\ref{sec:GW_linearizzata}), la loro emissione da parte di sistemi binari (Sezione~\ref{sec:GW_binarie}) e il fondo stocastico prodotto dalla popolazione cosmica di binarie di buchi neri supermassicci (Sezione~\ref{sec:GWB}), che è il segnale cercato dai \textit{Pulsar Timing Array}.

%%%%%%%%%%%%%%%%%%%%%%%%%%%%%%%%%%%%%%%%%%%%%%%%%%%%%%%%%%%%%%%%%%%%%%%%%%%%%%%%%%%%%%%%%%
\section{Onde gravitazionali in relatività generale linearizzata} \label{sec:GW_linearizzata}

In relatività generale la geometria dello spaziotempo, descritta dalla metrica $g_{\mu\nu}$, è legata al contenuto di materia ed energia dalle equazioni di campo di Einstein
\begin{equation}
    G_{\mu\nu} \equiv R_{\mu\nu} - \frac{1}{2} R\, g_{\mu\nu} = \frac{8\pi G}{c^4}\, T_{\mu\nu} \, ,
    \label{eq:einstein}
\end{equation}
dove $R_{\mu\nu}$ e $R$ sono il tensore e lo scalare di Ricci e $T_{\mu\nu}$ è il tensore energia-impulso. Queste equazioni sono non lineari, ma lontano dalle sorgenti il campo gravitazionale è debole e si può scrivere la metrica come una piccola perturbazione dello spaziotempo piatto di Minkowski~\cite{Maggiore2008}:
\begin{equation}
    g_{\mu\nu} = \eta_{\mu\nu} + h_{\mu\nu}\, , \qquad |h_{\mu\nu}| \ll 1 \, ,
    \label{eq:metrica_perturbata}
\end{equation}
con $\eta_{\mu\nu} = \mathrm{diag}(-1, 1, 1, 1)$. Sviluppando le equazioni~\eqref{eq:einstein} al primo ordine in $h_{\mu\nu}$ e introducendo la perturbazione a traccia invertita $\bar{h}_{\mu\nu} = h_{\mu\nu} - \frac{1}{2}\eta_{\mu\nu} h$ (con $h = \eta^{\mu\nu}h_{\mu\nu}$), si può sfruttare la libertà di scelta delle coordinate per imporre la condizione di gauge di Lorenz $\partial^{\nu}\bar{h}_{\mu\nu} = 0$. In questa gauge le equazioni di campo linearizzate diventano
\begin{equation}
    \Box\, \bar{h}_{\mu\nu} = -\frac{16\pi G}{c^4}\, T_{\mu\nu}\, ,
    \label{eq:einstein_linearizzata}
\end{equation}
dove $\Box = \eta^{\mu\nu}\partial_\mu\partial_\nu$ è l'operatore di d'Alembert. Nel vuoto ($T_{\mu\nu}=0$) la~\eqref{eq:einstein_linearizzata} si riduce a un'equazione delle onde, $\Box\,\bar{h}_{\mu\nu} = 0$: le perturbazioni della metrica si propagano come onde alla velocità della luce.

Non tutte le componenti di $\bar{h}_{\mu\nu}$ sono però fisiche. Dei 10 gradi di libertà di un tensore simmetrico $4\times4$, la gauge di Lorenz ne elimina 4; la libertà residua di trasformazione delle coordinate ne elimina altri 4. Si può allora scegliere la cosiddetta gauge \textit{transverse-traceless} (TT), definita dalle condizioni~\cite{Maggiore2008}
\begin{equation}
    h_{0\mu} = 0\, , \qquad h^{i}{}_{i} = 0 \, , \qquad \partial^{j} h_{ij} = 0 \, ,
    \label{eq:gauge_TT}
\end{equation}
in cui restano soltanto 2 gradi di libertà, che corrispondono alle due polarizzazioni dell'onda. Per un'onda piana monocromatica di pulsazione $\omega$ che si propaga lungo l'asse $z$ si ha
\begin{equation}
    h_{ij}^{\mathrm{TT}}(t,z) =
    \begin{pmatrix}
        h_+ & h_\times & 0 \\
        h_\times & -h_+ & 0 \\
        0 & 0 & 0
    \end{pmatrix}
    \cos\!\left[\omega\left(t - \frac{z}{c}\right)\right] ,
    \label{eq:onda_piana_TT}
\end{equation}
dove $h_+$ e $h_\times$ sono le ampiezze delle polarizzazioni ``più'' e ``per''. L'onda è quindi \textit{trasversale}: agisce solo nel piano perpendicolare alla direzione di propagazione.

L'effetto fisico di un'onda gravitazionale è una deformazione di marea delle distanze proprie. Due masse di prova libere poste a distanza $L$ lungo l'asse $x$ vedono la loro distanza propria oscillare di
\begin{equation}
    \frac{\delta L}{L} \simeq \frac{1}{2}\, h_{xx}^{\mathrm{TT}} = \frac{1}{2}\, h_+ \cos(\omega t)\, ,
    \label{eq:strain}
\end{equation}
motivo per cui l'ampiezza $h$ di un'onda gravitazionale è chiamata \textit{strain}. La polarizzazione $+$ allunga le distanze lungo $x$ mentre le contrae lungo $y$ (e viceversa mezzo periodo dopo); la polarizzazione $\times$ produce lo stesso effetto ruotato di $45^\circ$ (Figura~\ref{fig:polarizzazioni}).

\begin{figure}[h!]
    \centering
    % \includegraphics[width=0.8\textwidth]{Figure/cap1/polarizzazioni.pdf}
    \fbox{\parbox{0.8\textwidth}{\centering\vspace{1.5cm}\TODO{figura: anello di masse di prova deformato dalle polarizzazioni $+$ e $\times$}\vspace{1.5cm}}}
    \caption{\footnotesize Effetto delle due polarizzazioni di un'onda gravitazionale che si propaga lungo $z$ su un anello di masse di prova nel piano $xy$, nell'arco di un periodo.}
    \label{fig:polarizzazioni}
\end{figure}

%%%%%%%%%%%%%%%%%%%%%%%%%%%%%%%%%%%%%%%%%%%%%%%%%%%%%%%%%%%%%%%%%%%%%%%%%%%%%%%%%%%%%%%%%%
\section{Emissione di onde gravitazionali da sistemi binari} \label{sec:GW_binarie}

\subsection{La formula di quadrupolo}

Risolvendo la~\eqref{eq:einstein_linearizzata} lontano da una sorgente non relativistica, il termine dominante della radiazione emessa è proporzionale alla derivata seconda del momento di quadrupolo della distribuzione di massa~\cite{Maggiore2008}:
\begin{equation}
    h_{ij}^{\mathrm{TT}}(t, r) = \frac{2G}{c^4\, r}\, \Lambda_{ij,kl}\, \ddot{Q}_{kl}\!\left(t - \frac{r}{c}\right) ,
    \qquad
    Q_{ij} = \int \dd^3x\, \rho(t,\bm{x}) \left( x_i x_j - \frac{1}{3}r^2 \delta_{ij} \right) ,
    \label{eq:quadrupolo}
\end{equation}
dove $r$ è la distanza dalla sorgente, $\rho$ la densità di massa e $\Lambda_{ij,kl}$ il proiettore che seleziona la parte trasversa e a traccia nulla. A differenza dell'elettromagnetismo non esiste radiazione di monopolo né di dipolo: il monopolo di massa è la massa-energia totale, che si conserva, e il dipolo di massa è legato alla posizione del centro di massa, la cui derivata seconda è nulla per la conservazione della quantità di moto. La radiazione gravitazionale è quindi, al primo ordine, di natura quadrupolare. La potenza irraggiata è data dalla formula di quadrupolo di Einstein
\begin{equation}
    L_{\mathrm{GW}} = \frac{G}{5c^5} \left\langle \dddot{Q}_{ij}\, \dddot{Q}_{ij} \right\rangle ,
    \label{eq:luminosita_quadrupolo}
\end{equation}
dove $\langle\cdot\rangle$ indica una media su più periodi dell'onda.

\subsection{Binarie in orbita circolare}

Le sorgenti più importanti di onde gravitazionali sono i sistemi binari di oggetti compatti (nane bianche, stelle di neutroni, buchi neri). Per due masse $m_1$ e $m_2$ in orbita circolare con frequenza orbitale $f_{\mathrm{orb}}$, applicando la~\eqref{eq:quadrupolo} si trova che l'onda è emessa alla frequenza
\begin{equation}
    f = 2 f_{\mathrm{orb}} \, ,
\end{equation}
conseguenza della natura quadrupolare della radiazione (il momento di quadrupolo torna uguale a sé stesso dopo mezza orbita). L'ampiezza dell'onda, mediata sull'orientazione, dipende dalle masse solo attraverso la combinazione detta \textit{chirp mass}~\cite{Maggiore2008}:
\begin{equation}
    \mathcal{M} = \frac{(m_1 m_2)^{3/5}}{(m_1+m_2)^{1/5}}\, , \qquad
    h \simeq \frac{4}{r} \left(\frac{G\mathcal{M}}{c^2}\right)^{5/3} \left(\frac{\pi f}{c}\right)^{2/3} .
    \label{eq:ampiezza_binaria}
\end{equation}

L'energia emessa in onde gravitazionali viene sottratta all'energia orbitale: le due masse si avvicinano, la frequenza aumenta e con essa l'ampiezza. Bilanciando la~\eqref{eq:luminosita_quadrupolo} con la variazione dell'energia orbitale si ottiene l'evoluzione in frequenza
\begin{equation}
    \frac{\dd f}{\dd t} = \frac{96}{5}\, \pi^{8/3} \left(\frac{G\mathcal{M}}{c^3}\right)^{5/3} f^{11/3} \, ,
    \label{eq:chirp}
\end{equation}
il caratteristico \textit{chirp} che porta il sistema alla coalescenza. La~\eqref{eq:chirp} ha una conseguenza importante: poiché $\dd f/\dd t \propto f^{11/3}$, una binaria evolve molto lentamente a basse frequenze e trascorre la maggior parte della sua vita (nella banda di emissione) lontano dalla coalescenza.

\subsection{Binarie di buchi neri supermassicci}

La frequenza massima raggiungibile prima della coalescenza è, a meno di fattori di ordine unità, quella dell'ultima orbita circolare stabile (ISCO) di un buco nero di massa $M = m_1 + m_2$:
\begin{equation}
    f_{\mathrm{ISCO}} = \frac{c^3}{6^{3/2}\pi G M} \simeq 4.4\ \mathrm{kHz} \left(\frac{\Msun}{M}\right) .
    \label{eq:fisco}
\end{equation}
La frequenza tipica di emissione scala quindi come l'inverso della massa del sistema. Per binarie di massa stellare ($M \sim 10$--$100\,\Msun$) si trovano frequenze da decine a centinaia di Hz, la banda degli interferometri terrestri come LIGO e Virgo. Per le binarie di buchi neri supermassicci (\textit{supermassive black hole binaries}, SMBHB), con $M \sim 10^8$--$10^{10}\,\Msun$, la~\eqref{eq:fisco} dà frequenze al più dell'ordine del $\upmu$Hz, e durante la lunga fase di spiraleggiamento che precede la coalescenza (periodi orbitali di mesi o anni) l'emissione avviene nella banda dei nanohertz, $f \sim 10^{-9}$--$10^{-7}$~Hz~\cite{Taylor2021}. Questa è la banda di frequenze accessibile ai \textit{Pulsar Timing Array}, descritti nel Capitolo~\ref{chpt:PTA}.

Le SMBHB si formano in seguito alla fusione di due galassie, ciascuna delle quali ospita al proprio centro un buco nero supermassiccio: dopo la fusione i due buchi neri affondano verso il centro della nuova galassia per attrito dinamico e formano un sistema legato che, nelle fasi finali, perde energia principalmente per emissione di onde gravitazionali~\cite{Sesana2008, Taylor2021}.

\begin{figure}[h!]
    \centering
    % \includegraphics[width=0.85\textwidth]{Figure/cap1/spettro_GW.pdf}
    \fbox{\parbox{0.85\textwidth}{\centering\vspace{1.5cm}\TODO{figura: spettro delle onde gravitazionali (PTA / LISA / LIGO-Virgo) con le relative sorgenti --- da una review, citando la fonte}\vspace{1.5cm}}}
    \caption{\footnotesize Lo spettro delle onde gravitazionali: sorgenti e rivelatori nelle diverse bande di frequenza.}
    \label{fig:spettro_GW}
\end{figure}

%%%%%%%%%%%%%%%%%%%%%%%%%%%%%%%%%%%%%%%%%%%%%%%%%%%%%%%%%%%%%%%%%%%%%%%%%%%%%%%%%%%%%%%%%%
\section{Il fondo stocastico di onde gravitazionali} \label{sec:GWB}

Nell'Universo sono presenti moltissime SMBHB che emettono contemporaneamente nella banda dei nanohertz. Poiché, per la~\eqref{eq:chirp}, ciascuna di esse resta a lungo a bassa frequenza, il numero di sorgenti per intervallo di frequenza è elevato e i loro segnali si sommano in modo incoerente: invece di singole sorgenti risolte ci si aspetta di osservare un \textit{fondo stocastico di onde gravitazionali} (\textit{gravitational wave background}, GWB)~\cite{Sesana2008, Taylor2021}.

Un segnale stocastico non è prevedibile istante per istante e va descritto attraverso le sue proprietà statistiche. Assumendo che il fondo sia gaussiano, stazionario, isotropo e non polarizzato, tutta l'informazione è contenuta nella sua densità spettrale di potenza $S_h(f)$, o equivalentemente nella densità di energia per intervallo logaritmico di frequenza, normalizzata alla densità di energia critica dell'Universo $\varepsilon_c = 3H_0^2c^2/8\pi G$~\cite{Taylor2021}:
\begin{equation}
    \Omega_{\mathrm{GW}}(f) \equiv \frac{1}{\varepsilon_c} \frac{\dd \varepsilon_{\mathrm{GW}}}{\dd \ln f} = \frac{2\pi^2}{3H_0^2}\, f^3 S_h(f) = \frac{2\pi^2}{3H_0^2}\, f^2 h_c^2(f)\, ,
    \label{eq:omega_gw}
\end{equation}
dove si è introdotto lo \textit{strain caratteristico} $h_c(f) \equiv \sqrt{f\,S_h(f)}$.

Per una popolazione di binarie circolari la cui evoluzione è guidata dall'emissione di onde gravitazionali, lo spettro può essere calcolato con un argomento dovuto a Phinney~\cite{Phinney2001}: l'energia emessa da una singola binaria per intervallo logaritmico di frequenza scala come $\dd E/\dd\ln f \propto \mathcal{M}^{5/3} f^{2/3}$, quindi, indipendentemente dai dettagli della popolazione, $\Omega_{\mathrm{GW}}(f) \propto f^{2/3}$. Dalla~\eqref{eq:omega_gw} segue che lo strain caratteristico è una legge di potenza
\begin{equation}
    h_c(f) = \Agwb \left(\frac{f}{\fyr}\right)^{-2/3} ,
    \label{eq:hc_powerlaw}
\end{equation}
dove $\fyr = 1\ \mathrm{yr}^{-1} \simeq 3.17\times10^{-8}$~Hz è una frequenza di riferimento e $\Agwb$ l'ampiezza del fondo a quella frequenza. Tutta l'informazione astrofisica sulla popolazione di SMBHB (tasso di fusione delle galassie, distribuzione delle masse, \dots) è contenuta nell'ampiezza, i cui valori previsti sono dell'ordine di $10^{-15}$~\cite{Sesana2008}. Deviazioni dalla pendenza $-2/3$ alle frequenze più basse possono comparire se le binarie interagiscono ancora con l'ambiente (stelle, gas) o hanno orbite eccentriche~\cite{Taylor2021}.

Come vedremo nel Capitolo~\ref{chpt:PTA}, ciò che un PTA misura non è lo strain ma i ritardi nei tempi di arrivo degli impulsi delle pulsar. Il fondo descritto dalla~\eqref{eq:hc_powerlaw} produce in questi ritardi un segnale con densità spettrale di potenza~\cite{Taylor2021}
\begin{equation}
    S_{\mathrm{GWB}}(f) = \frac{h_c^2(f)}{12\pi^2 f^3} = \frac{\Agwb^2}{12\pi^2}\left(\frac{f}{\fyr}\right)^{-\gammagwb} \fyr^{-3}\, , \qquad \gammagwb = \frac{13}{3}\, ,
    \label{eq:psd_gwb}
\end{equation}
cioè ancora una legge di potenza, molto ripida, con più potenza alle basse frequenze: un segnale ``rosso''. I due parametri $(\Agwb, \gammagwb)$ sono quelli usati in tutta la tesi per descrivere il GWB. Come discusso nel Capitolo~\ref{chpt:rumore}, la stessa forma funzionale descrive anche diverse sorgenti di rumore delle pulsar: questa somiglianza è all'origine delle degenerazioni studiate in questo lavoro.
